%% =====================================================================
\section{Glauber 区}\label{sec:ir-glauber}

软与超软区的计数之所以简单，是因为程函近似在那里成立：软动量流过共线线时，分母只通过一个光锥分量感知它。有一个领头区域不满足这一点。

\begin{definition}[Glauber 区]\label{def:ir-glauber}
软动量 $k$ 的 \emph{Glauber 区}是 $|k^+k^-|\ll k_T^2$ 的区域，典型标度为 $Q(\lambda^2,\lambda^2,\lambda)$。在这一区域中，流过 $A$ 共线线 $p_A$ 的动量 $k$ 使
\begin{equation}\label{eq:ir-glauber-denom}
(p_A+k)^2\simeq p_A^2+2p_A^+k^--2\mathbf p_{A,T}\cdot\mathbf k_T-k_T^2 ,
\end{equation}
% VF-REG-7, VF-WL-11
其中 $k_T^2$ 与 $2p_A^+k^-$ 同阶，不能略去：只保留 $2p_A^+k^-$ 的程函近似的相对误差为 $O(1)$。Glauber 胶子的传播子 $\simeq-1/k_T^2$（瞬时的“Coulomb”交换），它以 $A^-$ 分量与 $A$ 共线线耦合（$p_A\cdot A\simeq p_A^+A^-$），以 $A^+$ 分量与 $B$ 共线线耦合。
\uses{def:ir-scalings}
\end{definition}

Glauber 区的幂次计数与软、超软区相同：在 $A$ 线与 $B$ 线之间交换一个胶子，Glauber 标度给出测度 $\lambda^6$、胶子传播子 $\lambda^{-2}$、两条线的分母各 $\lambda^{-2}$，合计 $\lambda^0$（软与超软标度分别给出 $\lambda^{4-2-1-1}$ 与 $\lambda^{8-4-2-2}$）。区别只在程函近似是否成立，所以关键在于 Glauber 区是否被钳住：若否，可把围道移到程函近似成立之处；若是，就必须另行处理。

\begin{proposition}[单圈 Glauber 区的钳夹判据 \tagR]\label{prop:ir-glauber-trapping}
（$R_{\mathrm{PT}}$：在重整化微扰论中对红外正规化的部分子态严格。）
(a)~单圈无质量顶点 $\gamma^*(q)\to q(p_1)\bar q(p_2)$（$p_1=(p_1^+,0,\mathbf 0_T)$，$p_2=(0,p_2^-,\mathbf 0_T)$，分母 $D_0=k^2=2k^+k^--k_T^2$、$D_1=(p_1-k)^2=-2(p_1^+-k^+)k^--k_T^2$、$D_2=(p_2+k)^2=2k^+(p_2^-+k^-)-k_T^2$，各带 $+\ii0$）在 Glauber 区没有钳夹点。具体地，固定 $k^+\sim\lambda^2Q$、$k_T\sim\lambda Q$ 时，三个分母的 $k^-$ 极点中只有 $D_1$ 的极点 $k^-\simeq-k_T^2/(2p_1^+)+\ii0$ 位于 $|k^-|\lesssim\lambda^2Q$ 的 Glauber 窗口内，$D_0$ 与 $D_2$ 的极点都在 $|k^-|\sim Q$ 处；对 $k^+$ 同理只有 $D_2$ 的一个极点。每个光锥积分在 Glauber 窗口内只遇到一个极点，围道可以移出该窗口，到 $|k^\pm|\sim k_T$ 处程函近似成立。
(b)~若同一共线子图中携带 $k$ 的既有一条入射的 $A$ 共线线 $p_A-k$，又有一条出射的 $A$ 共线线 $p_{A'}+k$（$\mathbf p_{A,T}=\mathbf p_{A',T}=0$，$p_A^-$、$p_{A'}^-$ 可略），其分母 $-2p_A^+k^--k_T^2+\ii0$ 与 $2p_{A'}^+k^--k_T^2+\ii0$ 的极点
\begin{equation}\label{eq:ir-glauber-poles}
k^-=-\frac{k_T^2}{2p_A^+}+\ii0,\qquad k^-=\frac{k_T^2}{2p_{A'}^+}-\ii0
\end{equation}
分居实轴两侧，在 $|k^-|\sim k_T^2/p^+\sim\lambda^2Q$ 处钳住 $k^-$ 围道。精确的 $k^-$ 积分有限，
\begin{equation}\label{eq:ir-glauber-exact}
\int_{-\infty}^{\infty}\frac{\dd k^-}{(-2p_A^+k^--k_T^2+\ii0)(2p_{A'}^+k^--k_T^2+\ii0)}=\frac{2\pi\ii}{2k_T^2\,(p_A^++p_{A'}^+)} ,
\end{equation}
% VF-REG-9
而略去 $k_T^2$ 的程函近似给出 $-\pi/[\delta(p_A^++p_{A'}^+)]$（$\ii0=\ii\delta$），在 $\delta\to0$ 时发散。（本命题只涉及 Feynman 积分。）
\uses{def:ir-glauber,thm:ir-landau-momentum}
\end{proposition}
\begin{proof}
(a)~$D_0=0$ 与 $D_2=0$ 的 $k^-$ 根 $(k_T^2-\ii0)/(2k^+)$、$-p_2^-+(k_T^2-\ii0)/(2k^+)$ 都是 $O(Q)$，$D_1=0$ 的根 $\simeq-k_T^2/(2p_1^+)+\ii0$ 为 $O(\lambda^2Q)$（数值拟合指数 $0$、$0$、$2$）。所以窗口内只有 $D_1$ 的极点，$k^-$ 围道可以向下半平面推离它直到 $|k^-|\sim k_T\ll Q$；对 $k^+$ 同理只有 $D_2$ 的极点。这与\cref{thm:ir-landau-momentum}一致：Glauber 区中 $D_0\simeq-k_T^2\ne0$，在壳线至多是 $D_1$、$D_2$，而 $u_1\simeq-2p_1$ 与 $u_2\simeq2p_2$ 不能以非负系数相消，没有 Landau 点。
(b)~被积函数在 $|k^-|\to\infty$ 时按 $(k^-)^{-2}$ 衰减，在上半平面闭合只包围第一个极点 $k^-_1=(-k_T^2+\ii0)/(2p_A^+)$。该处第一个因子对 $k^-$ 的导数为 $-2p_A^+$，第二个因子的值为 $2p_{A'}^+k^-_1-k_T^2=-k_T^2(p_A^++p_{A'}^+)/p_A^+$，故积分 $=2\pi\ii\cdot\frac{1}{-2p_A^+}\cdot\frac{-p_A^+}{k_T^2(p_A^++p_{A'}^+)}$，即\cref{eq:ir-glauber-exact}。略去 $k_T^2$ 后两个极点 $k^-=\ii\delta/(2p_A^+)$ 与 $k^-=-\ii\delta/(2p_{A'}^+)$ 间距为 $O(\delta)$，同样的留数计算给出 $2\pi\ii\cdot\frac1{-2p_A^+}\cdot\frac{1}{\ii\delta(p_A^++p_{A'}^+)/p_A^+}=-\frac{\pi}{\delta(p_A^++p_{A'}^+)}$。两式及 (a) 中的极点位置都已数值核对（相对误差 $<10^{-8}$）。
\end{proof}

(a) 与 CSS 综述对顶点 Glauber 区的分析一致：每个光锥变量只出现在一个分母里，积分不被钳住，可以形变到 $|k^\pm|\sim k_T$\cite[第~8.4~节，式~(137)--(138)]{CollinsSoperSterman1989}。注意两个分量都要形变：只把 $k^-$ 推到 $|k^-|\sim k_T$ 而保持 $k^+\sim\lambda^2Q$，仍有 $|k^+k^-|\sim\lambda^3Q^2\ll k_T^2$，并未离开 Glauber 区；%
% (misconception M31)
一般图中两个分量须沿 $e_G=n_B-n_A$ 协调形变（第~9~章的\cref{rem:method-two-component}与\cref{prop:method-glauber-deformation}）。(b) 是 Drell--Yan 中初态与末态旁观者相互作用的情形：同一共线子图中携带 Glauber 动量的线既有在硬散射之前的、又有在其后的，Glauber 区被钳住，程函近似不能使用，必须借助对末态割线的求和（幺正性）另证其相消——这在第~9、12~章与附录~O 中完成。在光锥序微扰论中，Glauber 钳夹还可以出现在不同项之间（项间钳夹），它阻碍 Ward 恒等式的直接使用，这也是 Glauber 区必须在割图求和的层次上处理的原因之一。
